\documentclass[tikz,border=10pt]{standalone}
\usepackage{tikz}
\usepackage{amsmath}
\usepackage{amssymb}
\usepackage{ctex}
\usetikzlibrary{calc, positioning, arrows.meta, matrix, decorations.pathreplacing, fit, backgrounds}

\begin{document}
\begin{tikzpicture}[>=Stealth,
  matbox/.style={draw=gray!60, fill=gray!5, rounded corners=4pt,
                 inner sep=8pt, align=center, font=\small},
]

% ================================================================
% 左侧：矩阵 A = [a1 a2 a3]
% ================================================================
\node[matbox, minimum width=2.8cm, minimum height=3.5cm]
  (boxA) at (0, 0) {};
\node[font=\small\bfseries, gray!60] at (0, 2.0) {$\mathbb{R}^m$ 中的列向量};

% 三列向量（用彩色小箭头示意，原始非正交）
\draw[->, thick, red!70!black] ($(boxA.center)+(-0.8,-0.8)$)
  -- ($(boxA.center)+(-0.4,0.8)$)
  node[above, font=\scriptsize] {$\mathbf{a}_1$};

\draw[->, thick, orange!70!black] ($(boxA.center)+(0,-0.9)$)
  -- ($(boxA.center)+(0.5,0.6)$)
  node[above, font=\scriptsize] {$\mathbf{a}_2$};

\draw[->, thick, purple!70!black] ($(boxA.center)+(0.8,-0.7)$)
  -- ($(boxA.center)+(0.3,0.9)$)
  node[above, font=\scriptsize] {$\mathbf{a}_3$};

% 标签
\node[font=\normalsize\bfseries] at (0, -2.1) {$A$};
\node[font=\scriptsize, gray!70] at (0, -2.5) {$m\times n$ 矩阵};

% ================================================================
% 等号
% ================================================================
\node[font=\Large\bfseries] at (2.5, 0) {$=$};

% ================================================================
% 中间：矩阵 Q = [q1 q2 q3]（正交）
% ================================================================
\node[matbox, minimum width=2.8cm, minimum height=3.5cm,
      draw=blue!50, fill=blue!5]
  (boxQ) at (5.0, 0) {};
\node[font=\small\bfseries, blue!60] at (5.0, 2.0) {标准正交列};

% 三个正交单位向量（两两垂直，用直角标记暗示）
\draw[->, thick, blue!80!black] ($(boxQ.center)+(-0.7,-0.5)$)
  -- ($(boxQ.center)+(-0.7,1.0)$)
  node[above, font=\scriptsize] {$\mathbf{q}_1$};

\draw[->, thick, teal!70!black] ($(boxQ.center)+(-0.05,-0.5)$)
  -- ($(boxQ.center)+(0.7,0.4)$)
  node[right, font=\scriptsize] {$\mathbf{q}_2$};

\draw[->, thick, green!60!black] ($(boxQ.center)+(0.7,-0.5)$)
  -- ($(boxQ.center)+(-0.2,0.8)$)
  node[above, font=\scriptsize] {$\mathbf{q}_3$};

% 小直角符号（q1 ⊥ q2 示意）
\draw[gray!50, thin]
  ($(boxQ.center)+(-0.7,0.1)+(0.1,0)$) --
  ($(boxQ.center)+(-0.7,0.1)+(0.1,0.1)$) --
  ($(boxQ.center)+(-0.7,0.1)+(0,0.1)$);

\node[font=\normalsize\bfseries, blue!80!black] at (5.0, -2.1) {$Q$};
\node[font=\scriptsize, blue!60] at (5.0, -2.5) {正交矩阵 $Q^TQ=I$};

% ================================================================
% 乘号
% ================================================================
\node[font=\Large\bfseries] at (7.5, 0) {$\cdot$};

% ================================================================
% 右侧：上三角矩阵 R
% ================================================================
\node[matbox, minimum width=3.2cm, minimum height=3.5cm,
      draw=orange!50, fill=orange!5]
  (boxR) at (10.0, 0) {};
\node[font=\small\bfseries, orange!60] at (10.0, 2.0) {上三角矩阵};

% 上三角示意（用数字格式展示）
\node[font=\small, align=center] at (10.0, 0.3) {
  $\begin{pmatrix}
    r_{11} & r_{12} & r_{13} \\
    0      & r_{22} & r_{23} \\
    0      & 0      & r_{33}
  \end{pmatrix}$
};

\node[font=\normalsize\bfseries, orange!80!black] at (10.0, -2.1) {$R$};
\node[font=\scriptsize, orange!60] at (10.0, -2.5) {$r_{ii}>0$（规范形式）};

% ================================================================
% 关系说明
% ================================================================
% Q^T Q = I 与 A=QR 的关键性质
\node[font=\small, fill=white, draw=gray!50, thin, rounded corners=3pt,
      inner sep=7pt, align=center]
  at (5.0, -3.8) {
    QR 分解：$A = QR$\\
    $Q^TQ = I$（正交矩阵），$R$ 为上三角（对角元为正）\\
    Gram-Schmidt 正交化 $\Leftrightarrow$ QR 分解
  };

% ================================================================
% Gram-Schmidt → QR 对应说明（上方括号）
% ================================================================
\draw[decorate, decoration={brace, amplitude=6pt, mirror},
      gray!60, thick]
  ($(boxA.north west)+(0,0.1)$) -- ($(boxA.north east)+(0,0.1)$)
  node[midway, above=8pt, font=\scriptsize, gray!70] {原始基};

\draw[decorate, decoration={brace, amplitude=6pt, mirror},
      blue!50, thick]
  ($(boxQ.north west)+(0,0.1)$) -- ($(boxQ.north east)+(0,0.1)$)
  node[midway, above=8pt, font=\scriptsize, blue!70] {正交基（单位化）};

% 整体标题
\node[font=\large\bfseries] at (5.0, 4.0) {QR 分解};

\end{tikzpicture}
\end{document}
